\documentclass[10pt,a4paper]{article}

\usepackage[margin=1.55cm,top=1.45cm,bottom=1.45cm]{geometry}
\usepackage[T1]{fontenc}
\usepackage{lmodern}
\usepackage{microtype}
\usepackage{xcolor}
\usepackage{enumitem}
\usepackage{hyperref}
\usepackage{titlesec}
\usepackage{tabularx}
\usepackage{array}

\definecolor{navy}{HTML}{17365D}
\definecolor{muted}{HTML}{4B5563}
\hypersetup{
  colorlinks=true,
  linkcolor=navy,
  urlcolor=navy,
  citecolor=navy,
  pdfauthor={Team 40},
  pdftitle={A Comprehensive Analysis of the Indian Startup Ecosystem}
}
\titleformat{\section}{\large\bfseries\color{navy}}{}{0em}{}[\titlerule]
\titleformat{\subsection}{\normalsize\bfseries\color{navy}}{}{0em}{}
\titlespacing*{\section}{0pt}{7pt}{4pt}
\titlespacing*{\subsection}{0pt}{5pt}{2pt}
\setlength{\parindent}{0pt}
\setlength{\parskip}{3pt}
\setlist[itemize]{leftmargin=*,nosep,topsep=2pt}
\setlist[enumerate]{leftmargin=*,nosep,topsep=2pt}
\urlstyle{same}
\sloppy

\begin{document}

\begin{center}
  {\LARGE\bfseries\color{navy} Team 40: A Comprehensive Analysis of the Indian Startup Ecosystem}\\[4pt]
  {\small
  B. Avanith (25B1013) \quad|\quad Shanmukha Sai (25B0907)\\
  Tejesh (25B0905) \quad|\quad Anirudh (25B1021)}
\end{center}
\vspace{-2pt}


\section{Project Hypothesis and Research Questions}

This project aims to investigate the evolution of India's startup
ecosystem between 2015 and 2021, focusing on changes in the geographical
and sectoral distribution of startup funding. It will then examine the
characteristics of startups that achieved unicorn status, exploring
whether common patterns emerge among these successful companies.

\subsection*{Primary Research Question}
\textbf{How did the city and sector composition of Indian startup
funding change from 2015 to 2021?}

We will investigate the following sub-questions:

\begin{enumerate}
    \item \textbf{Is Bangalore's dominance in India's startup ecosystem
    a thing of the past?} How has its share of startup funding changed
    relative to other major hubs such as Delhi-NCR and Mumbai, and
    have emerging cities gained ground?

    \item \textbf{Who's getting the money?} Which sectors have attracted
    the most startup funding, and have India's investment priorities
    shifted between 2015 and 2021?

    \item \textbf{Did COVID-19 rewrite the startup playbook?} How did
    the pandemic coincide with changes in funding volumes, sectoral
    preferences and the geographical distribution of startup funding?
\end{enumerate}

\subsection*{Follow-up Research Question}
\textbf{What does the path to unicorn status among Indian startups
look like?}

We will explore the following sub-questions:

\begin{enumerate}
    \item \textbf{What does the timeline look like?} How long do Indian
    startups typically take to reach a valuation of US\$1 billion,
    and how does it vary across companies?

    \item \textbf{Does location or industry determine a unicorn's
    destiny?} Which cities and sectors have produced the most
    Indian unicorns?

    \item \textbf{What's the unicorn formula?} What patterns emerge
    in unicorns' valuations and annual revenues?
\end{enumerate}

Using startup funding records and a separate dataset of Indian unicorns,
we will apply data cleaning, descriptive statistics and visualisation
techniques to document trends. The COVID-19 analysis and, similarly,
the unicorn analysis will be exploratory, looking for associated
changes (correlation) rather than implying causation. Differences in
coverage, missing values and other inconsistencies will be documented
when interpreting the findings.


\section{Datasets}
\begin{enumerate}
  \item \textbf{Startup Funding Dataset (2015--2021).} Monthly/tabular records of funding raised by Indian startups, with fields including city, founding date, industry, investors and amount (USD).\\
  \textit{Source:} \href{https://www.kaggle.com/datasets/arpan129/startups-funding-dataset}{Kaggle -- Startups Funding Dataset}.
  \item \textbf{List of 118 Indian Unicorn Startups (May 2025).} A curated list extracted from Tracxn, with location, sector, time-to-unicorn, valuation, investors and annual revenue.\\
  \textit{Source:} \href{https://www.kaggle.com/datasets/srisankargiri/list-of-118-unicorn-startups-in-indiamay-2025}{Kaggle -- Indian Unicorn Startups}.
\end{enumerate}

\section{Background Sources}
\begin{itemize}
  \item Growth91, \textit{Best Startup Sectors in India to Invest In for High Growth (2026)}: \href{https://growth91.com/blog/best-startup-sectors-india-invest-high-growth-2026/}{article link}.
  \item Saurabh Kapoor, \textit{Explained: What is a unicorn, and what does it take to become one?}, \textit{The Indian Express}: \href{https://indianexpress.com/article/explained/everyday-explainers/explained-what-is-a-unicorn-and-what-does-it-take-to-become-one-7944468/}{article link}.
  \item Think School, \textit{How Bangalore became the Silicon Valley of Asia?}: \href{https://youtu.be/cjhlcIOgUy4}{YouTube case study}.
  \item Nikhil Kamath, \textit{Inside India's Next Gen Startups}: \href{https://youtu.be/yFyIBKK361E}{YouTube podcast}.
\end{itemize}

\section{Intended Languages, Tools and Software}
We plan to use Python as the primary language, with pandas for loading, cleaning and grouping tabular data;
SciPy.stats for  statistical analysis ; and Matplotlib and Seaborn for plots. We will be using Plotly to create interactive charts and a lightweight Plotly
Dash dashboard to present city-wise, sector-wise and year-wise patterns. Jupyter Notebook will support exploration and GitHub will be used for version control and collaboration. We will consider R or
Tableau/Power BI only as alternatives if they offer a clear advantage for a particular analysis or dashboard.

\section{Logistics and Work Plan}
\textbf{Coordination and schedule:} The team will follow the project schedule and task outline recorded in the shared \href{https://docs.google.com/spreadsheets/d/1rOHg7PRTUyQhNcKS4sHRMWOS8t5IsXoPo_IaoyhDdhA/edit?usp=sharing}{logistics spreadsheet}, updating it as the project progresses. We will coordinate through GitHub and a WhatsApp group chat, and work in our hostel rooms, reading halls and CC503 as appropriate.

\textbf{Proposed division of work:}
\begin{itemize}
  \item \textbf{Brainstorming, data collection and data cleaning:} as a team.
  \item \textbf{Until the first checkpoint -- coding and analysis:} Shanmukha and Tejesh; \textbf{visualisations:} Avanith and Anirudh.
  \item \textbf{After the first checkpoint -- coding and analysis:} Avanith and Anirudh; \textbf{visualisations:} Shanmukha and Tejesh.
  \item \textbf{Status report:} Avanith. \quad
  \textbf{Video presentation:} Avanith and Shanmukha.
  \item \textbf{Video editing:} Tejesh. \quad
  \textbf{Final report:} Avanith.
\end{itemize}

\textbf{Help requested from the teaching team:} No specific assistance has been identified yet. We may seek guidance regarding the appropriate metrics for analysis,software use etc. as the project progresses.
\end{document}